% chapters/ch10-case-study-library.tex
\chapter{Case Study: Library System (End-to-End)}

\begin{learningobjectives}
In this chapter you will apply the full workflow to a library domain.
You will learn how to model titles versus physical copies, how to represent time-based events like loans, and how to encode key constraints such as “one active loan per copy”.
You will also practise validation with questions, traceability, and controlled test data.
\end{learningobjectives}

\section{Narrative}
The library domain is a strong teaching example because it forces two modelling decisions that appear in many real systems:
distinguishing an abstract item from its physical instances, and modelling events over time.

\begin{examplebox}
\textbf{Library narrative}\\
A library registers members.
The library stores book titles, and each title can have multiple physical copies.
Copies belong to a branch (a physical location).
Members borrow copies through loans, with checkout and due dates.
A copy may be returned later.
A copy cannot be on loan to two members at the same time.
Members can also place reservations for titles.
Reservations are handled in first-come-first-served order.
Overdue loans can generate fines, and fines can be paid.
\end{examplebox}

The phrase “title can have multiple copies” is the central constraint.
If you collapse title and copy, the model becomes confusing and will produce anomalies.

\section{Step 1: Business rules}
We write rules as testable statements.
We keep them focused on meaning and constraints.

\begin{examplebox}
\textbf{Business rules (library)}\\
R1: A member is uniquely identified by \texttt{member\_id}.\\
R2: A title is uniquely identified by \texttt{title\_id}.\\
R3: A copy is uniquely identified by \texttt{copy\_id}.\\
R4: Each copy belongs to exactly one title.\\
R5: A title can have zero or many copies.\\
R6: Each copy belongs to exactly one branch.\\
R7: A loan is for exactly one member and exactly one copy.\\
R8: A loan has a checkout date and a due date.\\
R9: A loan may have a return date (if returned).\\
R10: A copy has at most one active loan at a time.\\
R11: A reservation is made by exactly one member for exactly one title.\\
R12: Reservations are served in first-come-first-served order by request date.\\
R13: A member can have zero or many reservations over time.\\
R14: A title can have zero or many reservations over time.\\
R15: An overdue loan can generate a fine for the member.\\
R16: A fine can have zero or many payments.
\end{examplebox}

Rule R9 is the core business constraint.
It prevents an impossible state that would otherwise be representable.

\section{Step 2: Conceptual model}
We model the following entities:
Member, Title, Copy, Branch, Loan, Reservation, Fine, Payment.
Loan and Reservation are entities because they carry their own meaning and attributes.
They are not mere relationships.

\subsection{Conceptual decisions}
Write the meaning in sentences:
\begin{itemize}
  \item A title is an abstract bibliographic item; a copy is a physical instance.
  \item Copies belong to exactly one title; titles can have many copies.
  \item Copies belong to exactly one branch.
  \item Loans represent borrowing events over time and must be recorded historically.
  \item A reservation is for a title (not for a specific copy), because copies may not be available yet.
  \item Fines are linked to overdue loans; payments are linked to fines.
\end{itemize}

These decisions are what make the model coherent.

\subsection{ER-style diagram (minimal)}
\begin{center}
\resizebox{0.95\textwidth}{!}{%
\begin{tikzpicture}[node distance=18mm]
  \erentity{EMember}{}{Member}{member\_id}{name\\email}
  \erentity{ETitle}{right=55mm of EMember}{Title}{title\_id}{isbn\\title\\author}
  \erentity{EBranch}{right=55mm of ETitle}{Branch}{branch\_id}{name\\address}
  \erentity{ECopy}{below=25mm of ETitle}{Copy}{copy\_id}{barcode\\status}
  \erentity{ELoan}{below=25mm of EMember}{Loan}{loan\_id}{checkout\_date\\due\_date\\return\_date?}
  \erentity{ERes}{below=25mm of $(EMember)!0.5!(ETitle)$}{Reservation}{reservation\_id}{request\_date\\status}
  \erentity{EFine}{below=25mm of ELoan}{Fine}{fine\_id}{amount\\assessed\_date\\status}
  \erentity{EPay}{below=25mm of EFine}{Payment}{payment\_id}{paid\_at\\amount}

  \errel{RHas}{below=7mm of ETitle}{has}
  \erconnect{ETitle}{RHas}{1}{0..N}
  \erconnect{RHas}{ECopy}{1}{1}

  \errel{RLocated}{below=7mm of EBranch}{located\_at}
  \erconnect{EBranch}{RLocated}{1}{0..N}
  \erconnect{RLocated}{ECopy}{1}{1}

  % Loan edges: conceptual connections to Member and Copy
  \draw[er line] (ELoan) -- (EMember)
    node[pos=0.15,er card,above] {1}
    node[pos=0.85,er card,above] {1};
  \draw[er line] (ELoan) -- (ECopy)
    node[pos=0.15,er card,above] {1}
    node[pos=0.85,er card,above] {1};

  % Reservation edges: conceptual connections to Member and Title
  \draw[er line] (ERes) -- (EMember)
    node[pos=0.15,er card,above] {1}
    node[pos=0.85,er card,above] {1};
  \draw[er line] (ERes) -- (ETitle)
    node[pos=0.15,er card,above] {1}
    node[pos=0.85,er card,above] {1};

  % Fine edges: conceptual connections to Member and Loan
  \draw[er line] (EFine) -- (EMember)
    node[pos=0.15,er card,above] {1}
    node[pos=0.85,er card,above] {1};
  \draw[er line] (EFine) -- (ELoan)
    node[pos=0.15,er card,above] {1}
    node[pos=0.85,er card,above] {1};

  % Payment edges: conceptual connection to Fine
  \draw[er line] (EPay) -- (EFine)
    node[pos=0.15,er card,above] {0..N}
    node[pos=0.85,er card,above] {1};
\end{tikzpicture}%
}
\end{center}

\begin{figure}[ht]
\centering
\includegraphics[width=0.75\textwidth]{library-case.pdf}
\caption{The library diagram separates titles from copies and models loans and reservations as events.}
\end{figure}

The diagram is intentionally simple.
The “one active loan per copy” constraint is not drawn; it is specified as a constraint.

\section{Step 3: Relational schema (logical, no SQL)}
Now we translate meaning into relations with keys and references.

\[
\rel{Member}(\underline{member\_id}, name, email)
\]
\[
\rel{Title}(\underline{title\_id}, isbn, title, author)
\]
\[
\rel{Branch}(\underline{branch\_id}, name, address)
\]
\[
\rel{Copy}(\underline{copy\_id}, barcode, status, title\_id \rightarrow \rel{Title}, branch\_id \rightarrow \rel{Branch})
\]
\[
\begin{aligned}
\rel{Loan}(&\underline{loan\_id}, checkout\_date, due\_date, return\_date?, \\
&member\_id \rightarrow \rel{Member}, copy\_id \rightarrow \rel{Copy})
\end{aligned}
\]
\[
\begin{aligned}
\rel{Reservation}(&\underline{reservation\_id}, request\_date, status, \\
&member\_id \rightarrow \rel{Member}, title\_id \rightarrow \rel{Title})
\end{aligned}
\]
\[
\begin{aligned}
\rel{Fine}(&\underline{fine\_id}, amount, assessed\_date, status, \\
&member\_id \rightarrow \rel{Member}, loan\_id \rightarrow \rel{Loan})
\end{aligned}
\]
\[
\begin{aligned}
\rel{Payment}(&\underline{payment\_id}, paid\_at, amount, \\
&fine\_id \rightarrow \rel{Fine})
\end{aligned}
\]

\subsection{Attribute notes}
Two attributes deserve explicit meaning:
\begin{itemize}
  \item \texttt{Copy.status} expresses the library’s view of a copy (available, lost, repair, ...).
  \item \texttt{Reservation.status} expresses the state of a reservation (active, fulfilled, cancelled, ...).
  \item \texttt{Fine.status} expresses the state of a fine (open, paid, waived, ...).
\end{itemize}

In later chapters you may model these as reference entities.
For now, we keep them as constrained domains.

\subsection{Constraints (conceptual specification)}
We now specify the constraints that preserve meaning.

\begin{examplebox}
\textbf{Constraints (library)}\\
C1: Member.member\_id is unique and mandatory.\\
C2: Title.title\_id is unique and mandatory.\\
C3: Copy.copy\_id is unique and mandatory.\\
C4: Branch.branch\_id is unique and mandatory.\\
C5: Copy.title\_id references an existing Title and is mandatory.\\
C6: Copy.branch\_id references an existing Branch and is mandatory.\\
C7: Loan.member\_id references an existing Member and is mandatory.\\
C8: Loan.copy\_id references an existing Copy and is mandatory.\\
C9: Loan.due\_date is on or after Loan.checkout\_date.\\
C10: A Copy has at most one active Loan at a time.\\
\hspace*{1.6em}(Active loan = loan with return\_date missing, or status = active.)\\
C11: Reservation.member\_id references an existing Member and is mandatory.\\
C12: Reservation.title\_id references an existing Title and is mandatory.\\
C13: Reservations for a title are served by ascending request\_date.\\
C14: Fine.loan\_id references an existing Loan and is mandatory.\\
C15: Payment.fine\_id references an existing Fine and is mandatory.
\end{examplebox}

Constraint C10 is the key business constraint.
It is the one that protects the meaning “a copy cannot be loaned to two people at the same time”.
The definition of “active” must be explicit for this constraint to be enforceable or detectable.

In this model, a loan is active when \texttt{return\_date} is missing (or status = active, if you model status).
This definition should appear in the constraint catalog and in validation scenarios.

\section{Branches, fines, and queue policy}
These extensions change the operational meaning of the system.
They also introduce new constraints that must be explicit.

\subsection{Branches}
Branches are physical locations.
Copies belong to one branch, and availability is branch-specific.
This matters for questions like “which branch has an available copy?”

\subsection{Fines and payments}
Fines are linked to overdue loans and to the member responsible.
Payments are linked to fines and can be partial.
This keeps money flow separate from loan events while maintaining traceability.

\subsection{Reservation queue}
Reservations are served in first-come-first-served order by request date.
This is a policy rule, not a database default.
It must be written as a constraint and reflected in validation.

\begin{figure}[htbp]
\centering
\includegraphics[width=0.8\textwidth]{reservation-queue.pdf}
\caption{Reservation queues preserve first-come-first-served order for each title.}
\end{figure}

\section{Step 5: Validation without SQL}

\subsection{Key questions}
Choose questions that test the purpose of the system and the central constraints.

\begin{itemize}
  \item Q1: Which copies are available for a given title?
  \item Q2: What active loans does a member have?
  \item Q3: What is the loan history of a copy?
  \item Q4: Which loans are overdue today?
  \item Q5: Are there any copies with more than one active loan?
  \item Q6: What is the reservation queue for a title?
  \item Q7: Which branch holds the available copies of a title?
  \item Q8: What fines are outstanding for a member?
\end{itemize}

The schema supports these questions naturally:
availability comes from copies and active loans,
history comes from the loan records,
and reservations connect members to titles.

\subsection{Rule-to-schema traceability}
We link rules to design elements.

\begin{examplebox}
\textbf{Traceability map (library)}\\[4pt]
\begin{tabularx}{\textwidth}{@{}lX@{}}
\toprule
Rule & Design element(s) \\
\midrule
R4, R5 & Copy.title\_id FK to Title expresses 1:N (C5) \\
R6 & Copy.branch\_id FK to Branch expresses 1:N (C6) \\
R7 & Loan has mandatory FKs to Member and Copy (C7, C8) \\
R8, R9 & Loan has checkout/due dates and optional return\_date (C9) \\
R10 & One active loan per copy (C10) \\
R11 & Reservation has mandatory FKs to Member and Title (C11, C12) \\
R12 & Reservation queue order by request\_date (C13) \\
R13, R14 & Multiple reservations possible because no uniqueness on those FKs \\
R15 & Fine references Loan and Member (C14) \\
R16 & Payment references Fine (C15) \\
\bottomrule
\end{tabularx}
\end{examplebox}

\subsection{Controlled test data (paper test)}
Construct these scenarios:
\begin{itemize}
  \item A title with two copies, one available and one on loan.
  \item A copy loaned and returned, creating history.
  \item Two loans for the same copy with missing return\_date (invalid under C10).
  \item Two reservations for the same title (valid).
\end{itemize}

If the invalid state is not clearly representable as a violation, you need stronger constraint specification.
If you cannot represent history cleanly, you likely collapsed events into static relationships.

\subsection{Paper dataset: valid scenario}
\begin{examplebox}
\textbf{Valid dataset (excerpt)}\par\smallskip
\begin{tabular}{@{}llll@{}}
\toprule
copy\_id & title\_id & branch\_id & status \\
\midrule
C1 & T1 & B1 & available \\
C2 & T1 & B1 & on\_loan \\
\bottomrule
\end{tabular}
\par\smallskip
\begin{tabular}{@{}lllll@{}}
\toprule
loan\_id & copy\_id & member\_id & checkout\_date & return\_date \\
\midrule
L1 & C2 & M1 & 2026-02-01 & \text{missing} \\
L2 & C2 & M1 & 2025-12-10 & 2025-12-20 \\
\bottomrule
\end{tabular}
\end{examplebox}

\subsection{Paper dataset: invalid scenario}
\begin{examplebox}
\textbf{Invalid dataset (violates C10)}\par\smallskip
\begin{tabular}{@{}lllll@{}}
\toprule
loan\_id & copy\_id & member\_id & checkout\_date & return\_date \\
\midrule
L3 & C1 & M2 & 2026-02-03 & \text{missing} \\
L4 & C1 & M3 & 2026-02-04 & \text{missing} \\
\bottomrule
\end{tabular}
\end{examplebox}
Both loans are active for the same copy, violating the “one active loan per copy” rule.

\section{Design note: why reservations are for titles}
Some readers try to reserve a copy.
That seems concrete, but it is usually wrong at the conceptual level.
A reservation is a promise about future availability, and the library typically fulfils it with any available copy of the title.

Modelling reservations for titles also matches real operations:
a title can have multiple copies, and copies can be lost, repaired, or replaced.
If reservations are tied to a specific copy, the system becomes brittle.

This is a general modelling lesson:
when a concept is a policy about a class of items, it should reference the abstract entity, not a specific instance.

\section{Common mistakes}
Library models often fail in familiar ways.

Mistake: Collapsing Title and Copy.\\
Fix: Separate entities so availability and history are clean.

Mistake: Treating reservations as copy-level.\\
Fix: Reserve the title and allocate a copy when available.

Mistake: Leaving “active loan” undefined.\\
Fix: Define active explicitly (e.g., return\_date missing) and enforce it.

Mistake: Mixing fines into the loan table.\\
Fix: Keep fines as a separate entity linked to loans and members.

\section{Mini checklist}
Before finalizing the library model, check:
\begin{itemize}[leftmargin=*,nosep]
  \item Are Title and Copy separated with a clear 1:N relationship?
  \item Is “active loan” defined and used consistently?
  \item Does each copy belong to one branch?
  \item Are reservation queues ordered by request date?
  \item Are fines and payments modelled as separate entities?
  \item Can the schema represent loan history without duplication?
  \item Do constraints cover availability and queue policy?
\end{itemize}

\begin{chaptersummary}
This case study demonstrates the full workflow in a library domain.
The key modelling decision is separating titles from copies, and representing borrowing as a time-based event via loans.
Constraints, especially “one active loan per copy”, preserve correctness.
Validation through questions, traceability, and paper tests makes the design defensible before implementation.
\end{chaptersummary}

\begin{exercisebox}
Extend the library model.
Answer in full sentences, then update the schema in book notation.
\begin{enumerate}[label=\arabic*.]
  \item Add support for multiple authors per title. Decide whether Author is an entity and model it.
  \item Add fines for overdue loans. Decide what facts should be stored versus derived.
  \item Add a branch concept (multiple library locations) and decide where copies belong.
  \item Add a rule limiting maximum active loans per member and describe how it would be enforced or detected.
  \item Update the traceability map for your new rules and constraints.
  \item Define “active loan” precisely and write the corresponding constraint in words.
  \item Create a valid and an invalid paper dataset for reservations and explain the queue policy.
  \item Add payments for fines and list two constraints that protect correctness.
\end{enumerate}
\end{exercisebox}

In the next chapter we focus on state and time.
Loan history, reservation queues, and fine status will connect directly to those concepts.
